\documentclass[verbale=interno]{gpt}

\titolo{Discussione e votazione dei capitolati}
\data{2026-10-13}
\orario{17:40}{19:00}
\luogo{Discord}
\redattore{mazzurana}
\verificatore{cescobolla}
\componenti{biasini,cescobolla,corro,guglielmo,mazzurana,spagnoli}
\assenti{}

\modifica{1.0.0}{2026-10-15}{mazzurana}{cescobolla}{Verifica e approvazione}
\modifica{0.1.0}{2026-10-14}{mazzurana}{--}{Prima stesura}

\begin{document}

\section{Ordine del giorno}
\begin{enumerate}
  \item Presentazione dei capitolati di interesse da parte di ciascun membro.
  \item Selezione dei capitolati preferiti dal gruppo.
  \item Votazione finale tra i capitolati selezionati.
  \item Organizzazione del lavoro per la settimana successiva.
\end{enumerate}

\section{Svolgimento}

\subsection{Presentazione dei capitolati}
Ogni membro ha presentato al gruppo i capitolati che aveva approfondito
dopo l'incontro precedente, illustrandone i punti di interesse. Ne è
seguita una discussione comune.

\subsection{Prima votazione}
Al termine della discussione il gruppo ha votato i capitolati preferiti.
I più votati sono risultati:
\begin{itemize}
  \item \textbf{C0} -- AIRA (Sanmarco Informatica S.p.A.);
  \item \textbf{C1} -- AURA (Vimar S.p.A.);
  \item \textbf{C2} -- NEXUM Concierge (Eggon S.r.l.).
\end{itemize}

\subsection{Votazione finale}
Il gruppo ha poi votato il capitolato preferito tra i soli tre
selezionati. Ogni membro ha espresso un voto; i risultati sono riportati
in Tabella~\ref{tab:voti}.

\begin{table}[H]
  \centering
  \caption{Risultati della votazione finale\label{tab:voti}}
  \begin{tabular}{@{}lllrr@{}}
    \toprule
    \leavevmode{\sffamily\scriptsize\bfseries\color{gptgrigio}POSIZIONE} &
    \leavevmode{\sffamily\scriptsize\bfseries\color{gptgrigio}CAPITOLATO} &
    \leavevmode{\sffamily\scriptsize\bfseries\color{gptgrigio}PROPONENTE} &
    \leavevmode{\sffamily\scriptsize\bfseries\color{gptgrigio}VOTI} &
    \leavevmode{\sffamily\scriptsize\bfseries\color{gptgrigio}PERCENTUALE} \\
    \midrule
    1\textsuperscript{o} & \textbf{C2} -- NEXUM Concierge & Eggon S.r.l.
      & 3 & 50\% \\
    2\textsuperscript{o} & \textbf{C1} -- AURA & Vimar S.p.A.
      & 2 & 33\% \\
    3\textsuperscript{o} & \textbf{C0} -- AIRA & Sanmarco Informatica S.p.A.
      & 1 & 17\% \\
    \bottomrule
  \end{tabular}
\end{table}

\subsection{Organizzazione del lavoro}
Per approfondire i tre capitolati il gruppo si è diviso in coppie. Ogni
coppia raccoglie informazioni sul capitolato assegnato e redige una
lista di possibili domande da porre al proponente. Le domande saranno
riviste e validate insieme nel prossimo incontro e, se approvate,
inviate alle aziende.

\section{Decisioni}
\begin{decisioni}
  \decisione{I capitolati di interesse del gruppo sono C0, C1 e C2.}
  \decisione{L'ordine di preferenza è: 1)~C2 -- NEXUM Concierge,
    2)~C1 -- AURA, 3)~C0 -- AIRA.}
  \decisione{L'approfondimento dei capitolati avviene a coppie:
    C0 a Luca Cesco Bolla e Alessio Corrò, C1 a Enrico Biasini e
    Francesco Guglielmo, C2 a Francesco Mazzurana e Gaetano Spagnoli.}
  \decisione{Le domande per i proponenti vengono validate da tutto il
    gruppo prima dell'invio.}
  \decisione{I verbali vengono redatti con un template comune.}
\end{decisioni}

\section{Azioni da intraprendere}
\begin{azioni}
  \azione{Preparare il template per i verbali.}{guglielmo}{19/10/2026}
  \azione{Raccogliere informazioni su C0 -- AIRA e preparare le
    domande per il proponente.}{cescobolla,corro}{19/10/2026}
  \azione{Raccogliere informazioni su C1 -- AURA e preparare le
    domande per il proponente.}{biasini,guglielmo}{19/10/2026}
  \azione{Raccogliere informazioni su C2 -- NEXUM Concierge e
    preparare le domande per il proponente.}{mazzurana,spagnoli}{19/10/2026}
\end{azioni}

\section{Prossima riunione}
Lunedì 19 ottobre 2026 alle 21:00, per rivedere e validare le domande
preparate dalle coppie ed eventualmente inviarle alle aziende.

\end{document}
